\subsection{A sharp unforced profile obstruction above slope one half}
\label{unforced:obstruction-section}

The following example establishes sharpness only within the elementary
profile problem. It is not a planar set with a null distance set and does
not establish a necessary Hausdorff--packing condition.

\begin{proposition}[Exact single-collapse cost]\label{unforced:obstruction}
For \(0<a<b\le1\), let
$$
 g(x)=\min\{bx,1+a-x\}\qquad(0\le x\le1),
 \qquad C=\frac{b-a}{1+2b}.
$$
This is a 1-Lipschitz profile with \(ax\le g(x)\le bx\). For
\(0<\zeta<C\), the infimum of the costs of all admissible chains from
\(1-\zeta\) to zero, with no forced intermediate endpoint, is exactly
\(C-\zeta\). An attaining chain has
\(O(1+\log(1/\zeta))\) edges.
\end{proposition}

\begin{proof}
The two affine pieces meet at
$$
 p=\frac{1+a}{1+b}.
$$
Their slopes are \(b\) and \(-1\), so the Lipschitz assertion holds.
The lower barrier on the terminal piece follows from
\(1+a-x-ax=(1+a)(1-x)\ge0\), and all other barrier comparisons are
immediate. Put
$$
 s=1-C=\frac{1+a+b}{1+2b},\qquad
 q=2s-1=\frac{1+2a}{1+2b}.
$$
The assumption \(a<b\) gives \(q<p<s\).

Consider any chain, and its first edge from \(n>s\) to an endpoint
\(u\le s\). Every preceding edge lies on the terminal affine piece
of slope \(-1\), so the preceding costs sum to \(1-\zeta-n\).
If \(u\ge p\), the crossing edge costs \(n-u\ge n-s\). If
\(u<p\), admissibility gives \(u\ge2n-1>q\). Since the profile has
only one peak, its minimum over the crossing interval is at an endpoint,
and the crossing cost is
$$
 \bigl(bu-(1+a-n)\bigr)_+
 \ge(1+2b)(n-s).
$$
In either case, the total cost through this crossing edge is at least
\(C-\zeta\). All subsequent costs are nonnegative, proving the lower
bound.

For attainment, double the complementary depth from \(1-\zeta\)
down to \(s\), shortening the last jump if needed. This portion lies
on the terminal descending piece and costs exactly \(C-\zeta\).
The next edge \(s\to q\) is admissible with equality and has zero cost,
because
$$
 g(s)=1+a-s=\frac{b(1+2a)}{1+2b}=g(q),
$$
and the profile is above this common value between the endpoints. On
\([0,q]\), the profile is increasing, so all downward edges cost zero.
Maximal admissible jumps finish at zero without further cost. Their
number, together with the initial doubling part, has the stated
logarithmic bound.
\end{proof}

For \(b\ge1/2\), the limit \(C\) in this example matches the universal
upper bound in Lemma~\ref{unforced:profile}. Thus the limiting universal
profile constant is sharp in that range. The example remains valid for
\(b<1/2\), where equality with the universal minimax is not asserted.
When \(a=b\), the profile is simply \(ax\) and all costs are zero.

At \(a=1/4\), \(b=1\), the example has peak \(p=5/8\), critical
endpoints \(s=3/4\), \(q=1/2\), and limiting cost \(1/4\). It
therefore reaches the strict source-margin boundary at the full-packing
dimension pair \(d=5/4\), \(D=2\). This obstructs a uniform improvement
of this particular profile constant. It does not rule out a different
analytic argument, additional geometric information, or a different
finite-scale balance, and it supplies no distance-set counterexample.
